\begin{tikzpicture}[
    >=Stealth,
    input/.style={
        draw=black!75,
        line width=0.8pt,
        rounded corners=2pt,
        fill=gray!10,
        align=center,
        minimum height=0.95cm,
        text width=2.25cm,
        font=\small
    },
    proc/.style={
        draw=black!80,
        line width=0.8pt,
        rounded corners=2pt,
        fill=gray!5,
        align=center,
        minimum height=0.95cm,
        text width=2.5cm,
        font=\small
    },
    highlight/.style={
        draw=UdeAGreen!85!black,
        line width=0.8pt,
        rounded corners=2pt,
        fill=UdeAGreen!7,
        align=center,
        minimum height=0.95cm,
        text width=2.5cm,
        font=\small
    },
    outnode/.style={
        draw=black!80,
        line width=0.8pt,
        rounded corners=2pt,
        fill=white,
        align=center,
        minimum height=0.85cm,
        text width=2.35cm,
        font=\small
    },
    arrow/.style={->, line width=0.8pt, draw=black!75}
]

% ==============================================================================
% 1. ENTRADA Y FILTRADO PREVIO
% ==============================================================================
\node[input] (signal) {
    Señal adquirida\\
    $v_i[n]$
};

\node[proc, right=0.65cm of signal] (filter) {
    Filtrado pasabanda\\
    Banda de interés\\
    $y_i[n]$
};

\draw[arrow] (signal) -- (filter);

% ==============================================================================
% 2. BIFURCACIÓN CON ESPACIO HORIZONTAL HOLGADO
% ==============================================================================
% Tronco a 0.55 cm del filtro
\coordinate (fork) at ($(filter.east) + (0.55, 0.0)$);

% Bloques desplazados a 1.15 cm a la derecha del filtro (0.60 cm a la derecha del tronco)
\node[highlight, above right=0.70cm and 1.15cm of filter] (hilbert) {
    Transformada\\
    de Hilbert\\
    $z_i[n] = y_i[n] + j\widetilde{y}_i[n]$
};

\node[proc, below right=0.85cm and 1.15cm of filter] (gate) {
    Selección temporal\\
    Ventana directa\\
    $T_{\mathrm{gate}} = K / f_0$
};

\node[highlight, below=0.65cm of gate] (toa_proc) {
    Detección de arribo\\
    Umbral $V_{\mathrm{th}}$ e\\
    interpolación de cruce
};

% Trazado de flechas holgadas (0.60 cm de recorrido horizontal antes del recuadro)
\draw[line width=0.8pt, draw=black!75] (filter.east) -- (fork);
\draw[arrow] (fork) |- (hilbert.west);
\draw[arrow] (fork) |- (gate.west);
\draw[arrow] (fork) |- (toa_proc.west);

% ==============================================================================
% 3. BLOQUES DE SALIDA (ESTRUCTURA ORIGINAL DESPEJADA)
% ==============================================================================

% Salidas de Hilbert
\node[outnode, right=0.70cm of hilbert] (env) {
    Envolvente\\
    $A_i[n] = |z_i[n]|$
};

\node[outnode, above=0.25cm of env] (phase) {
    Fase instantánea\\
    $\Phi_i[n] = \operatorname{arg} z_i[n]$
};

\node[outnode, below=0.25cm of env] (freq) {
    Frecuencia inst.\\
    $f_{i,\mathrm{inst}}[n]$ (Savitzky-G.)
};

\draw[arrow] (hilbert.east) -- ++(0.30, 0.0) |- (phase.west);
\draw[arrow] (hilbert.east) -- (env.west);
\draw[arrow] (hilbert.east) -- ++(0.30, 0.0) |- (freq.west);

% Salida de energía (alineada a Selección temporal)
\node[outnode, right=0.70cm of gate] (energy) {
    Energía y RMS\\
    $Q_i = \sum (y_i - \bar{y})^2$\\
    $A_{\mathrm{RMS},i}$ (Síncrona)
};
\draw[arrow] (gate) -- (energy);

% Salida de ToA (alineada a Detección de arribo)
\node[outnode, right=0.70cm of toa_proc] (toa) {
    Tiempo de llegada\\
    $t_{\mathrm{arr},i}$
};
\draw[arrow] (toa_proc) -- (toa);

\end{tikzpicture}